\chapter{Tables and rows}\label{chap:rows}

This chapter describes how \ratpoints{} computes the tables of its moduli and how it builds and keeps
the rows that the sieve reads (Chapter~\ref{chap:sieve}). Chapter~\ref{chap:moduli} says what the
tables say and why their \emph{admission tests} let the numerator of every point through: the row
of a denominator~$b$ in the sieve entry of a modulus admits a numerator~$a$ when the tables allow
$F(a, b)$ to be a square. Each builder is proved to build the \emph{closed form} of its row, the
array of words whose bits are the values of the admission test, so that every row of a run lets
every point through. The contracts of the tables, which reason modulo~$p$ and~$p^e$, use Mathlib,
the others Lean core and \verb|Std| only; the proofs use that the constants $\mathrm{rbaPack}$ and
$\mathrm{chunk}$ (Chapter~\ref{chap:sieve}) are positive, never their values.

\section{The tables as \ratpoints{} computes them}\label{sec:rowsResidues}

\verb|examine_prime| (\csource{find_points.c}{1683}{1826}) computes the table of a prime~$p$
(Definition~\ref{def:tableOfPrime}) by Horner's rule in \verb|unsigned long|s with Barrett's
reduction, on a fixed schedule (\csource{find_points.c}{1628}{1655},
\csource{find_points.c}{1716}{1747}): below $\mathrm{maxPrimeEven} = 2^{\mathrm{maxBitsInPrime}}$
an accumulator below~$p$ survives $\mathrm{hornerSteps} = 64/\mathrm{maxBitsInPrime} - 1$ steps
$s \mapsto s a + c$ within 64 bits, and \ratpoints{} examines no larger prime. For a factor $q = p^e$
with $e \ge 2$, \verb|examine_power| (\csource{find_points.c}{2254}{2331}) computes the tables of
Definition~\ref{def:sieveEntry}.

\begin{lemma}[The table of a prime as \ratpoints{} computes it]\label{lem:hornerBarrett}
  \uses{def:tableOfPrime}
  \lean{Ratpoints.Sieve.examinePrime?, Ratpoints.Sieve.examinePrime?_spec, Ratpoints.hornerSteps, Ratpoints.Sieve.barrett, Ratpoints.Sieve.hornerBarrett, Ratpoints.Sieve.barrett_eq, Ratpoints.Sieve.HornerInvariant, Ratpoints.pow_hornerSteps_add_one_lt, Ratpoints.Sieve.hornerBarrett_spec}
  \leanok
  Barrett's reduction of~$u$ modulo~$p$ with the reciprocal~$m$ is
  $r = u - \lfloor u m / 2^{64} \rfloor\, p$, and $r - p$ when $r \ge p$
  (\csource{find_points.c}{1639}{1650}); the value of~$f$ at~$a$ modulo~$p$ is Horner's rule
  from~$c_d$ on the coefficients reduced modulo~$p$, with Barrett's reduction at the end and, for
  $d > \mathrm{hornerSteps}$, after every $\mathrm{hornerSteps}$ steps
  (\csource{find_points.c}{1726}{1731}, \csource{find_points.c}{1737}{1744}); the table is made of
  these values and handed back with the coefficients modulo~$p$, or none, the exit, when no entry
  of it is true (\csource{find_points.c}{1683}{1697}). Then:
  \begin{enumerate}
    \item for $p > 0$ and $u < 2^{64}$, Barrett's reduction with
      $m = \lfloor (2^{64} - 1)/p \rfloor$ gives $u \bmod p$;
    \item for $p < \mathrm{maxPrimeEven}$, $a < p$ and coefficients below~$p$, the schedule gives
      $f(a) \bmod p$, and the accumulator stays below $p^{\mathrm{hornerSteps} + 1} < 2^{64}$;
    \item for $0 < p < \mathrm{maxPrimeEven}$ the computed entry, the prime with its table and its
      inverses, is that of Definition~\ref{def:tableOfPrime}, with the coefficients modulo~$p$, or
      none exactly when no entry of that table is true.
  \end{enumerate}
\end{lemma}

\begin{proof}
  \uses{def:tableOfPrime}
  \leanok
  Write $2^{64} - 1 = m p + t$ with $t < p$; then $u/p - u m/2^{64} = u (1 + t)/(p\, 2^{64})$ lies
  in $[0, 1)$, so the quotient is $\lfloor u/p \rfloor$ or one less. The invariant of Horner's loop:
  the accumulator is congruent modulo~$p$ to the value of the polynomial of the coefficients read,
  and below $p^{k + 1}$, where $k \le \mathrm{hornerSteps}$ counts the steps since the last
  reduction, and $p^{64/\mathrm{maxBitsInPrime}} \le 2^{64}$ since
  $p < 2^{\mathrm{maxBitsInPrime}}$, whatever the degree. For~(3), both ways of reducing the
  coefficients agree, and the loop over the residues makes the entries of the reference table
  by~(2).
\end{proof}

\begin{lemma}[The tables of a prime power]\label{lem:modinv}
  \uses{def:sieveEntry}
  \lean{Ratpoints.Sieve.examinePower, Ratpoints.Sieve.examinePower_spec, Ratpoints.Sieve.modinv, Ratpoints.Sieve.ModinvInvariant, Ratpoints.Sieve.modinv_spec, Ratpoints.Sieve.unitInverses, Ratpoints.Sieve.unitInverses_spec, Ratpoints.modWords, Ratpoints.Sieve.setBitAt, Ratpoints.Sieve.setSquareBits, Ratpoints.Sieve.setSquareBits_spec, Ratpoints.maxPrimeEven_le_sixty_four_mul_modWords}
  \leanok
  The tables of $m = p^e$ as \verb|examine_power| makes them from the coefficients reduced
  modulo~$m$ (\csource{find_points.c}{2267}{2282}), in bit arrays of
  $\mathrm{modWords} = \lceil \mathrm{maxPrimeEven}/64 \rceil$ words, bit~$x$ in bit $x \bmod 64$
  of word $\lfloor x/64 \rfloor$ (\csource{rp-private.h}{523}{523}): the bits of the squares
  $x^2 \bmod m$, $x < m$, set in a bit array~\verb|sq| (\csource{find_points.c}{2283}{2284}); the
  bits of \verb|fsq| and \verb|gsq| set where the values of~$f$ and of the reversed polynomial, by
  Horner's rule, are marked in~\verb|sq|; and the inverses by the extended Euclidean algorithm
  \verb|modinv| (\csource{find_points.c}{2241}{2252}), $0$ for the residues divisible by~$p$
  (\csource{find_points.c}{2302}{2303}). Then:
  \begin{enumerate}
    \item for $1 < m$ and~$x$ prime to~$m$, \verb|modinv| gives the inverse of~$x$ modulo~$m$ in
      $[0, m)$; for a prime~$p$ and $e \ge 1$, every entry of the inverses is below~$p^e$, and the
      entry of a unit $x < p^e$ is its inverse;
    \item $\mathrm{maxPrimeEven} \le 64\,\mathrm{modWords}$, and for $m \le 64\,\mathrm{modWords}$
      a bit array of $\mathrm{modWords}$ words gets the bit of every square modulo~$m$ set;
    \item for a prime~$p$, $e \ge 1$ and $p^e \le 64\,\mathrm{modWords}$, the tables of~$p^e$ have
      these inverses, bit $x < p^e$ of \verb|fsq| set whenever~$f(x)$ is a square modulo~$p^e$, and
      bit $t < p^e$ of \verb|gsq|, for $p \mid t$, set whenever the reversed polynomial of the even
      degree $D = d + (d \bmod 2)$ takes a square value at~$t$; only this direction is needed, since
      a further bit loses no point.
  \end{enumerate}
\end{lemma}

\begin{proof}
  \leanok
  The Euclidean loop runs on remainders~$r_i$ and cofactors~$s_i$ from $r_0 = m$, $r_1 = x \bmod m$,
  $s_0 = 0$, $s_1 = 1$, with the invariant $r_i \equiv s_i x \bmod m$,
  $\gcd(r_0, r_1) = \gcd(m, x)$, $r_0 > 0$, cofactors of opposite signs with
  $|s_0|\, r_1 + |s_1|\, r_0 = m$; it ends since~$r_1$ decreases, with $r_0 = 1$, and keeps the
  cofactors in $[-m, m]$, the range of their \verb|long|s. For~(3), a value that is a square
  modulo~$m$ is $y^2 \bmod m$ for some $y < m$, so that its bit in~\verb|sq| is set by~(2), and the
  loops over the residues and the multiples of~$p$ set the bit of the residue.
\end{proof}

\section{The closed form of a row}\label{sec:closedRow}

A row of the modulus~$m$ has $m + \mathrm{chunk} - 1$ bit arrays of $\mathrm{rbaPack}$ words;
bit~$t$ of word~$j$ stands for the number $64 j + t$ (\csource{init.c}{252}{258}), and the last
$\mathrm{chunk} - 1$ bit arrays are copies of the first ones, for the chunked first stage. Which
numerator a bit stands for is the business of the packing (Chapter~\ref{chap:sieve}).

\begin{lemma}[The closed form of a row]\label{lem:closedRow}
  \uses{def:sieveEntry, def:firstStage}
  \lean{Ratpoints.Sieve.closedRow, Ratpoints.Sieve.wordOfBits, Ratpoints.Sieve.getLsbD_closedRow, Ratpoints.Sieve.getElem!_closedRow_mod, Ratpoints.Sieve.SieveEntry.wraps_closedRow, Ratpoints.Sieve.PatternBits, Ratpoints.Sieve.RowWords, Ratpoints.Sieve.RowWords.toArray_eq_closedRow}
  \leanok
  The \emph{closed form} of the row of a residue~$r$ modulo~$m$, for an admission test
  $\mathrm{adm}$, which says whether the row of~$r$ admits a number~$N$, is the array of
  $(m + \mathrm{chunk} - 1)\,\mathrm{rbaPack}$ words whose word~$j$ has bit~$t$ set exactly when
  $t < 64$ and $\mathrm{adm}(r)$ admits $64 j + t$. With the admission test of a sieve entry, of
  the table of a prime or of the tables of a prime power, it is the row that \ratpoints{} builds; it
  has no C code of its own. Then:
  \begin{enumerate}
    \item bit~$t$ of word $j < (m + \mathrm{chunk} - 1)\,\mathrm{rbaPack}$ is $t < 64$ and
      $\mathrm{adm}(r)(64 j + t)$;
    \item for $m > 0$ and a test that does not change when a multiple of~$m$ is added to the
      number, word~$j$ is word $j \bmod m\,\mathrm{rbaPack}$, so that the closed form of a row of a
      sieve entry of positive modulus repeats its first~$m$ bit arrays after them;
    \item a word \emph{holds} $n \le 64$ bits of a pattern~$g$ of period~$m$ from~$s$ on when its
      bit $t < n$ is $g((s + t) \bmod m)$ and its other bits are clear; if
      $g(N \bmod m) = \mathrm{adm}(r)(N)$ for every~$N$, an array of
      $(m + \mathrm{chunk} - 1)\,\mathrm{rbaPack}$ words whose word~$j$ holds the 64 bits of~$g$
      from~$64 j$ on is the closed form; the contracts of the builders conclude by this.
  \end{enumerate}
\end{lemma}

\begin{proof}
  \leanok
  Two words are equal when their bits are. For~(2), the numbers $64 j + t$ and
  $64 (j \bmod m\,\mathrm{rbaPack}) + t$ differ by a multiple of~$m$, and the test of a sieve entry
  reduces the numerator modulo its modulus first. For~(3), both bits are $g((64 j + t) \bmod m)$.
\end{proof}

\section{The tables of rows and the builders}\label{sec:rowBuilders}

A sieve entry of \ratpoints{} (\csource{rp-private.h}{542}{555}) holds the array \verb|sieve| of the
rows of its residues, \verb|NULL| for a row not built, and the pointer \verb|init| to its builder.
A prime has its row of~$0$ from the table \verb|sieves0| that the build generates
(\csource{gen_find_points_h.c}{102}{134}, \csource{find_points.c}{1820}{1821}) and the others
absent; a prime power has one entry, made when a modulus first needs it, and a product an entry of
its own that points to the entries of its factors (\csource{find_points.c}{2414}{2475}). The
builders are generated for each prime from \verb|CODE_INIT_SIEVE1| for $p < 64$
(\csource{init.c}{118}{183}) and from \verb|CODE_INIT_SIEVE2| for $p > 64$
(\csource{init.c}{186}{244}); a prime power has \verb|_ratpoints_sieve_init_power|
(\csource{init.c}{355}{393}) and a product \verb|_ratpoints_sieve_init_product|
(\csource{init.c}{395}{445}). A builder publishes its row in the entry, where it stays for the run.

\begin{definition}[The tables of rows]\label{def:rowStore}
  \lean{Ratpoints.Sieve.EntryRef, Ratpoints.Sieve.RowStore, Ratpoints.Sieve.RowStore.row?, Ratpoints.Sieve.RowStore.publish, Ratpoints.Sieve.RowInit, Ratpoints.Sieve.sieveInit}
  \leanok
  A \emph{reference} to an entry, for a pointer to it: a prime by itself, a prime power by its
  prime and its exponent, a product by the position of its modulus among the moduli. The
  \emph{store} maps each reference to its table of rows, $m$ rows for the modulus~$m$, each present
  or absent; \emph{publishing} a row puts it into the table of a reference and leaves a store
  without that table as it is. The kinds of builder, for the pointer \verb|init|
  (\csource{rp-private.h}{497}{498}): the function of a prime below~$64$ from
  \verb|CODE_INIT_SIEVE1| and that of a prime above~$64$ from \verb|CODE_INIT_SIEVE2|, which the
  primes of the table have (\csource{gen_init_sieve_h.c}{48}{62}), and the builders of a prime power
  and of a product. A sieve entry holds the reference to its table of rows and the kind of its
  builder: for a prime or a prime power those of its one factor, for a product its own
  (\lname{Ratpoints.Sieve.SieveEntry.refInit}).
\end{definition}

\begin{lemma}[The rows of a prime]\label{lem:primeRows}
  \uses{lem:closedRow, def:rowStore, def:tableOfPrime}
  \lean{Ratpoints.Sieve.sieves0, Ratpoints.Sieve.sieves0_eq_closedRow, Ratpoints.Sieve.extract_sieves0, Ratpoints.Sieve.initFourLoop, Ratpoints.Sieve.initFourLoop_spec, Ratpoints.Sieve.initReplicate, Ratpoints.Sieve.initReplicate_spec, Ratpoints.Sieve.sieveInit1, Ratpoints.Sieve.sieveInit1_spec, Ratpoints.Sieve.sieveInit2, Ratpoints.Sieve.sieveInit2_spec}
  \leanok
  The row of~$0$ of a prime~$p$: the~$p$ words of the forbidden divisor~$p$, whose word~$r$ has
  bit~$t$ clear exactly when $p \mid 64 r + t$ (\lname{Ratpoints.Sieve.forbiddenRow}), written
  $\mathrm{rbaPack}$ times and followed by the copies (\csource{gen_find_points_h.c}{122}{134}).
  The other rows by one function for the primes below~$64$ and one for those above, from the table
  of the prime, the reference, the residue~$b$ and the store: bit $a < p$ of the pattern the entry
  of the table at $a d \bmod p$ for the inverse~$d$ of~$b$, by a loop that steps four residues by
  $4 d \bmod p$ (\csource{init.c}{72}{83}), rotated into the first~$p$ words, the copies
  (\csource{init.c}{99}{111}), and the row published in the table of the reference
  (\csource{init.c}{178}{183}, \csource{init.c}{239}{244}). Then:
  \begin{enumerate}
    \item for $p > 0$ the row of~$0$ is the closed form of the row of~$0$ of every table of~$p$,
      and its first~$p$ words are those that the forbidden divisors read
      (Chapter~\ref{chap:denominators});
    \item for the table of a prime $p < 64$ and $0 < b < p$ whose inverse in the table is
      below~$p$, the function from \verb|CODE_INIT_SIEVE1| returns the closed form of the row
      of~$b$, and the store with it published in the table of the reference; the same holds for
      $p \ge 64$ with $p \bmod 64 \ne 0$ and the function from \verb|CODE_INIT_SIEVE2|.
  \end{enumerate}
  The builders need no primality and nothing of the table of squares.
\end{lemma}

\begin{proof}
  \uses{lem:closedRow}
  \leanok
  The loop keeps four lanes at the residues $s + (a + k)\, d \bmod p$, $k < 4$, and its word holds
  the bits below~$a$ of the pattern. \verb|CODE_INIT_SIEVE1| rotates the words of the row out of
  two words by the shifts $\delta = (\lfloor 64/p \rfloor + 1)\, p \bmod 64$ and $64 - \delta$,
  \verb|CODE_INIT_SIEVE2| out of a ring of $\lfloor p/64 \rfloor + 1$ words by $p \bmod 64$ and
  $64 - (p \bmod 64)$; every shift executed lies in $[0, 63]$, where the shifts of the C code are
  defined, since $\delta \in [1, 63]$ for $p < 64$ and $p \bmod 64 \ne 0$ otherwise. With the
  copies, word~$j$ holds the 64 bits from~$64 j$ on of the pattern read at $N d \bmod p$, the
  admission test at $b \ne 0$, and Lemma~\ref{lem:closedRow}(3) concludes. At $b = 0$ every lane
  would stay at~$0$, which is why the row of~$0$ comes from \verb|sieves0|, whose words are right
  since the multiples~$p i$, $i < 64$, are the multiples of~$p$ below~$64 p$.
\end{proof}

\begin{lemma}[The rows of a prime power]\label{lem:powerRows}
  \uses{lem:closedRow, def:rowStore, def:sieveEntry, lem:modinv}
  \lean{Ratpoints.Sieve.layOutPattern, Ratpoints.Sieve.PatClear, Ratpoints.Sieve.layOutPattern_spec, Ratpoints.Sieve.powerPatternUnit, Ratpoints.Sieve.powerPatternNonunit, Ratpoints.Sieve.sieveInitPower, Ratpoints.Sieve.sieveInitPower_spec}
  \leanok
  The layout of a pattern of period~$m$: word $x < m\,\mathrm{rbaPack}$ of the row the 64 bits from
  the offset $64 x \bmod m$ on of the pattern repeated through $w = \lfloor (m + 127)/64 \rfloor$
  words, then the copies (\csource{init.c}{259}{267}, \csource{init.c}{280}{297}). The row of the
  residue~$b$ of a prime power~$m$, from its tables, the reference and the store: for a unit~$b$,
  bit~$x$ of the pattern set when bit $x b^{-1} \bmod m$ of \verb|fsq| is
  (\csource{init.c}{369}{375}), and for $p \mid b$, bit~$x$ of a unit~$x$ set when bit
  $b x^{-1} \bmod m$ of \verb|gsq| is (\csource{init.c}{378}{385}); the pattern laid out, and the
  row published in the table of the reference (\csource{init.c}{391}{393}). Then:
  \begin{enumerate}
    \item for $0 < m \le 64\,\mathrm{modWords}$ and a pattern whose bits~$x$ with
      $m \le x < 64 \lceil m/64 \rceil$ are clear, word~$j$ of the row holds the 64 bits of the
      pattern, read with period~$m$, from~$64 j$ on;
    \item for $b < m \le 64\,\mathrm{modWords}$ whose entry in the table of inverses is below~$m$,
      the builder returns the closed form of the row of~$b$, and the store with it published in the
      table of the reference.
  \end{enumerate}
\end{lemma}

\begin{proof}
  \uses{lem:closedRow}
  \leanok
  The repeated pattern has at every bit $N < 64 w$ the bit $N \bmod m$ of the pattern, and
  $64 w \ge m + 64$, so the window at $64 x \bmod m$ is the window at~$64 x$ of the periodic
  pattern; the C code's guard against a shift by~64 keeps every shift in $[1, 63]$. The patterns
  set the bits that the admission test of the prime power reads: \verb|fsq| at~$x b^{-1}$ for a
  unit~$b$, and for $p \mid b$, \verb|gsq| at~$b x^{-1}$ for a unit~$x$ and nothing for $p \mid x$.
\end{proof}

\begin{lemma}[The rows of a product]\label{lem:productRows}
  \uses{lem:closedRow, def:rowStore, def:sieveEntry, lem:primeRows, lem:powerRows}
  \lean{Ratpoints.Sieve.initFactorRow, Ratpoints.Sieve.Factor.Buildable, Ratpoints.Sieve.Factor.Builds, Ratpoints.Sieve.initFactorRow_spec, Ratpoints.Sieve.factorRows, Ratpoints.Sieve.productCopy, Ratpoints.Sieve.productAnd, Ratpoints.Sieve.sieveInitProduct, Ratpoints.Sieve.FactorRowsPre, Ratpoints.Sieve.ProductRowsPre, Ratpoints.Sieve.RowStore.Grows, Ratpoints.Sieve.sieveInitProduct_spec}
  \leanok
  The row of the residue~$b$ of a product: the rows of its factors at~$b$ modulo their moduli, read
  from their tables or built into them through the pointer of the entry of each factor
  (\csource{init.c}{407}{414}); the first laid down in blocks of its length
  (\csource{init.c}{419}{423}), the others and-ed in (\csource{init.c}{424}{429}), the copies
  (\csource{init.c}{430}{430}), and the row published (\csource{init.c}{443}{445}). A factor is
  \emph{buildable} when its builder has what it needs: a prime $p > 0$ with its inverses below~$p$
  and $p \bmod 64 \ne 0$ for $p \ge 64$, or a prime power $0 < m \le 64\,\mathrm{modWords}$ with
  its inverses below~$m$; it may be called below its modulus, for a prime not at~$0$. Then:
  \begin{enumerate}
    \item a buildable factor called where it may be called returns the closed form of its row and
      publishes it in the table given;
    \item if the product has a positive modulus~$m$ and a factor, the moduli of its factors
      divide~$m$, its factors are buildable, the rows present in their tables are closed forms, a
      row absent below the modulus of a factor is one where its builder may be called, and two
      factors with one reference have the same closed forms, then the builder returns the closed
      form of the row of~$b$, and the store has \emph{grown} by rows of the factors (the tables
      keep their sizes, no row is replaced, every new row is the closed form of a factor whose
      table it is in), with the row of the product published.
  \end{enumerate}
\end{lemma}

\begin{proof}
  \uses{lem:closedRow}
  \leanok
  (1) is Lemma~\ref{lem:primeRows} or~\ref{lem:powerRows}, by the kind. For~(2), the invariant of
  the loop over the factors: their rows before the position of the loop are closed forms, and the
  store has grown by rows of the factors; a row read is a closed form also when a factor before it
  with the same reference built it. Each modulus~$q$ of a factor divides~$m$, so word~$y$ of the
  row is the and of the words $y \bmod q\,\mathrm{rbaPack}$ of the rows of the factors, as the
  admission test of the product is the and of theirs
  (\lname{Ratpoints.Sieve.toArray\_eq\_closedRow\_of\_productWord}).
\end{proof}

\section{The rows of a run}\label{sec:runRows}

The sieve sets up a denominator by \verb|set_rows| (\csource{find_points.c}{3149}{3170}), the row
of the residue of each modulus of the first two stages from its entry, perhaps absent. It builds
each absent row by the function of the entry, under the lock of the tables when threads share them
and after reading the pointer again, since another thread may have built it in the meantime
(\csource{find_points.c}{3522}{3554}); the model, which has no threads, reads it again as well.
Within a run the moduli of the first two stages are pairwise coprime (Chapter~\ref{chap:moduli}),
so that each table of rows is reached through one modulus; the contracts also admit entries that
share a table. The run hands the store on from one denominator and one block to the next.

\begin{definition}[The tables of rows of a run]\label{def:runTables}
  \uses{def:rowStore, def:sieveEntry, def:firstStage, lem:primeRows, lem:productRows, def:packing}
  \lean{Ratpoints.Sieve.primeRows, Ratpoints.Sieve.SieveEntry.addTables, Ratpoints.Sieve.SieveEntry.fetchRow, Ratpoints.Sieve.setRows, Ratpoints.Sieve.buildRows}
  \leanok
  The table of rows of a prime~$p$: its row of~$0$ that of Lemma~\ref{lem:primeRows}, the other
  $p - 1$ rows absent (\csource{find_points.c}{1820}{1821}). The tables of rows of a sieve entry,
  each made only when its reference is absent: that of each factor, and for a product its own
  (\csource{find_points.c}{2427}{2427}, \csource{find_points.c}{2471}{2471}). The row of a
  residue~$b$ of a sieve entry \emph{read or built}: the row in its table when it is there, and
  otherwise the row that its builder builds into it. The sieve specs of a denominator, as
  \verb|set_rows| makes them (\csource{find_points.c}{3149}{3170}): for each entry of the first two
  stages its modulus, the offset of the row of the class and its row at its residue of~$b$, absent
  when it is not built, with the end~$p$. The missing rows of a denominator: each sieve spec of the
  first two stages without a row gets the row read or built of its entry at its residue
  (\csource{find_points.c}{3522}{3554}).
\end{definition}

\begin{lemma}[The rows of a run are closed forms]\label{lem:runRows}
  \uses{def:runTables, lem:closedRow, def:sieveEntry, lem:sieveEntry_sound, def:firstStage, def:packing}
  \lean{Ratpoints.Sieve.SieveEntry.fetchRow_spec, Ratpoints.Sieve.StoreHolds, Ratpoints.Sieve.SieveEntry.RowsFit, Ratpoints.Sieve.RowStore.Fits, Ratpoints.Sieve.runRows, Ratpoints.Sieve.RowStore.MadeFor, Ratpoints.Sieve.RowStore.madeFor_nil, Ratpoints.Sieve.RowStore.MadeFor.foldl_addTables, Ratpoints.Sieve.RowStore.MadeFor.storeHolds, Ratpoints.IsModulusEntry.rowsFit_runRows, Ratpoints.Sieve.buildRows_setRows_map_spec, Ratpoints.Sieve.RowStore.Fits.buildRows_setRows}
  \leanok
  A store \emph{holds} a function~$R$ of the references and the residues for a list of entries when
  every row present in the table of a reference~$\rho$ at~$r$ is $R(\rho, r)$ and the table of every
  prime factor of the entries has its row of~$0$. An entry at the position~$n$ \emph{fits}~$R$ when
  it has the reference and the kind of builder of that position (and, with one factor, the modulus
  of that factor), a positive modulus, a factor, moduli of the factors dividing its modulus and
  buildable factors, and~$R$ at its reference and at those of its factors is the closed form of
  their rows; a store \emph{fits} a list of entries when it holds some~$R$ for them that every entry
  fits. The \emph{rows of the references} of a run are the closed forms of the rows of the table of
  a prime, of the tables of a prime power and of the entry of a product at its position; a store is
  \emph{made for} a list of entries when it holds the rows of the references of every run whose
  entries begin with the list, and every table of rows of a prime in it, or of a prime factor of the
  entries, has its row of~$0$. Then:
  \begin{enumerate}
    \item if the row present in the table of an entry at~$b$ is the closed form, and the builder
      has what it needs when the row is absent, the row read or built is the closed form of the
      row of~$b$, and the store is the same or has grown by rows of the factors, with the row
      published;
    \item the store of the loop of the tables, which holds the rows of~$0$ of the primes and no
      others, is made for no entries; adding the tables of rows of further entries, each with the
      reference and the kind of builder of its position and with positive primes, gives a store
      made for the longer list; a store made for the entries of a run holds the rows of their
      references for them; and an entry of a modulus as the choice makes it
      (Chapter~\ref{chap:moduli}) fits the rows of the references of the run at its position;
    \item for a store that holds~$R$ for a list~$L$, entries of~$L$ that fit~$R$ at their positions
      and residues below their moduli, the sieve specs that \verb|set_rows| makes and the missing
      rows complete have the closed forms of the rows of their entries at their residues, the
      offsets of the class, the start~$0$ and the end~$p$, and the store after them holds~$R$
      for~$L$; so a store that fits the entries of a run fits them after each denominator.
  \end{enumerate}
\end{lemma}

\begin{proof}
  \uses{lem:primeRows, lem:productRows, lem:modinv}
  \leanok
  (1) by the cases of the definition and Lemma~\ref{lem:productRows}. (2) The rows of~$0$ are
  \verb|sieves0|, the closed form for every table of the prime, and a table made only when absent
  holds no other row; two entries that reach one table have one closed form, since~$R$ depends on
  the reference only and the rows of a prime power do not depend on the switch of its density
  (\lname{Ratpoints.Sieve.admits\_run\_examinePower}). (3) The invariant of the loop over the
  specs: the store holds~$R$, and the specs before the position of the loop have their closed
  forms; a row present is~$R$ at its reference, a row built is a closed form by~(1), and a builder
  of a prime is never called at~$0$, whose row is present from the start.
\end{proof}

The choice of the moduli applies item~(2) to its entries
(\lname{Ratpoints.IsChoiceEntries.fits\_foldl\_addTables}), which are sound
(Lemma~\ref{lem:sieveEntry_sound}), and item~(3) hides the store from the sieve of one denominator
(Chapter~\ref{chap:sieve}).

\begin{lemma}[Every row lets every point through]\label{lem:rowsAdmitPoints}
  \uses{lem:sieveEntry_sound, def:sieveEntry, def:binaryForm, def:bruteForce}
  \lean{Ratpoints.Sieve.SieveEntry.IsSound.admits_of_sq_eq, Ratpoints.Sieve.survives_of_mem_pointsOfDen}
  \leanok
  The row of~$b$ of a sound entry admits~$a$ for every triple $(a, y, b)$ with $\gcd(a, b) = 1$
  and $y^2 = F(a, b)$, and the numerator of every point with the denominator~$b$ that the
  brute-force reference lists (\lname{Ratpoints.BruteForce.pointsOfDen}) survives every list of
  sound entries for~$b$.
\end{lemma}

\begin{proof}
  \uses{lem:sieveEntry_sound, def:bruteForce}
  \leanok
  A point of the reference has $\gcd(a, b) = 1$ and $y^2 = F(a, b)$, and a numerator survives a
  list of entries when every entry admits it.
\end{proof}

Every row that the run reads is the closed form of the admission test of its entry
(Lemma~\ref{lem:runRows}), whose bit at a number is the test at that number
(Lemma~\ref{lem:closedRow}); so the sieve never clears the bit of the numerator of a point, which
the membership theorem (Theorem~\ref{thm:membership}) needs, through the contract of the sieve of
one denominator (Definition~\ref{def:sift}) and the block lemma (Theorem~\ref{thm:blocks}).

\section{Where the model differs from the C code}\label{sec:rowsDeviations}

Chapter~\ref{chap:ccode} summarizes the general differences (the generated tables as functions of
the prime, two builder functions for the primes, the store for the arrays reached through pointers,
no threads). The \verb|long|s are natural numbers or integers, with the range obligations as
lemmas: the accumulator of Horner's rule below~$2^{64}$, the cofactors of the extended Euclidean
algorithm in $[-m, m]$, and the inverses of a prime power below~$2^{16}$, the range of their
\verb|unsigned short|
(\lname{Ratpoints.Sieve.getElem!\_inv\_run\_examinePower\_lt\_two\_pow\_16}). Barrett's reduction
has the form with the 128-bit product, which the default build compiles, and the branch of
the schedule of Horner's rule is chosen for each residue, where \verb|examine_prime| chooses it
once. The probes \verb|RP_INIT_NOACC| and \verb|RP_INIT_NOREP| (\csource{init.c}{88}{93},
\csource{init.c}{97}{98}), whose tables are wrong by design, and the check
\verb|RP_VERIFY_MODULI| (\csource{init.c}{299}{353}) are excluded.
